\chapter{Gasto público: los agregados}

\marginal{el ajuste recae en el programable}
El gasto neto total aprobado para 2025 es de \textbf{9,302,015.8 millones de
pesos}, la misma cifra que el total de ingresos, como manda la identidad. En
términos del PIB baja de 27.0 a \textbf{25.5\%}, y el gasto programable de 19.7 a
\textbf{17.8}. \textbf{Casi todo el ajuste de dos puntos del PIB recae sobre el
gasto programable}, que cae 7.3\% real contra el cierre estimado de 2024 y 4.1\%
contra el aprobado.

Esto no es una sorpresa: es la ejecución de un ajuste que el paquete anterior ya
había anunciado. El escenario de mediano plazo del CGPE 2024 proyectaba que el
gasto programable pagado bajaría de 18.8 a 16.4\% del PIB en un solo año. El
paquete 2025 lo entrega, con una magnitud algo menor.

\begin{table}[htbp]
\centering
\caption{Estimación de las finanzas públicas, 2024--2025}
\label{tab:C4_1}
\small
\begin{tabular}{lSSSSS}
\toprule
{Concepto} & {2024 (mdp)} & {2025 (mdp)} & {2024 \% PIB} & {2025 \% PIB} & {Var. real \%} \\
\midrule
{RFSP} & \num{-1864872.3} & \num{-1428348.1} & \num{-5.4} & \num{-3.9} & \num{-26.5} \\
{Recursos financieros extrapresupuestarios} & \num{-171872.3} & \num{-257781.7} & \num{-0.5} & \num{-0.7} & \num{43.9} \\
{Balance presupuestario} & \num{-1693000.0} & \num{-1170566.5} & \num{-4.9} & \num{-3.2} & \num{-33.7} \\
{Ingresos presupuestarios} & \num{7328995.2} & \num{8055649.4} & \num{21.3} & \num{22.3} & \num{5.4} \\
{Ingresos petroleros} & \num{1048069.4} & \num{1142021.5} & \num{3.0} & \num{3.2} & \num{4.5} \\
{Ingresos no petroleros} & \num{6280925.8} & \num{6913627.9} & \num{18.3} & \num{19.1} & \num{5.6} \\
{Gobierno Federal (no petroleros)} & \num{5203187.3} & \num{5670839.1} & \num{15.1} & \num{15.7} & \num{4.5} \\
{Tributarios} & \num{4941540.8} & \num{5296426.4} & \num{14.4} & \num{14.6} & \num{2.8} \\
{No tributarios} & \num{261646.5} & \num{374412.7} & \num{0.8} & \num{1.0} & \num{37.3} \\
{Organismos y empresas} & \num{1077738.5} & \num{1242788.8} & \num{3.1} & \num{3.4} & \num{10.6} \\
{Gasto neto pagado} & \num{9021995.2} & \num{9226215.8} & \num{26.2} & \num{25.5} & \num{-1.9} \\
{Gasto programable pagado} & \num{6451161.3} & \num{6451831.3} & \num{18.8} & \num{17.8} & \num{-4.1} \\
{Diferimiento de pagos} & \num{-44050.6} & \num{-75800.0} & \num{-0.1} & \num{-0.2} & \num{65.1} \\
{Gasto programable devengado} & \num{6495211.9} & \num{6527631.3} & \num{18.9} & \num{18.0} & \num{-3.6} \\
{Gasto no programable} & \num{2570833.8} & \num{2774384.5} & \num{7.5} & \num{7.7} & \num{3.5} \\
{Costo financiero} & \num{1263994.1} & \num{1388373.6} & \num{3.7} & \num{3.8} & \num{5.4} \\
{Participaciones} & \num{1262789.1} & \num{1340210.9} & \num{3.7} & \num{3.7} & \num{1.8} \\
{Adefas} & \num{44050.6} & \num{45800.0} & \num{0.1} & \num{0.1} & \num{-0.3} \\
{Balance primario presupuestario} & \num{-429005.9} & \num{217807.2} & \num{-1.2} & \num{0.6} & {---} \\
{SHRFSP} & \num{16787906.1} & \num{18591005.5} & \num{48.8} & \num{51.4} & \num{6.2} \\
\bottomrule
\end{tabular}
\par\vspace{2pt}\footnotesize Fuente: elaboración propia del ITED con información de la SHCP, CGPE 2025, Anexo II.6, p.\textasciitilde{}82. El CGPE 2025 se publicó sin capa de texto; las cifras se leyeron de la página renderizada y se validaron por identidad contable.
\end{table}


\section{Qué se mueve, por ramo}

\begin{table}[htbp]
\centering
\caption{Gasto por ramo del Gobierno Federal, aprobado contra aprobado (mdp)}
\label{tab:C4_2}
\footnotesize
\setlength{\tabcolsep}{4pt}
\begin{tabular}{L{4.2cm}SSSSS}
\toprule
{Concepto} & {2024} & {2025} & {Dif. nominal} & {Dif. real} & {Var. real \%} \\
\midrule
{19 Aportaciones a Seguridad Social} & \num{1421823.7} & \num{1480920.4} & \num{59096.7} & \num{-2041.7} & \num{-0.1} \\
{28 Participaciones a Entidades Federativas y Municipios} & \num{1262789.1} & \num{1340210.9} & \num{77421.8} & \num{23121.9} & \num{1.8} \\
{24 Deuda Pública} & \num{1023011.5} & \num{1149551.6} & \num{126540.1} & \num{82550.6} & \num{7.7} \\
{33 Aportaciones Federales para Entidades Federativas y Municipios} & \num{984484.5} & \num{979951.8} & \num{-4532.7} & \num{-46865.5} & \num{-4.6} \\
{20 Bienestar} & \num{543933.0} & \num{579883.9} & \num{35950.9} & \num{12561.7} & \num{2.2} \\
{11 Educación Pública} & \num{439017.9} & \num{465871.9} & \num{26853.9} & \num{7976.2} & \num{1.7} \\
{47 Entidades no Sectorizadas} & \num{137996.6} & \num{173656.4} & \num{35659.8} & \num{29725.9} & \num{20.7} \\
{23 Provisiones Salariales y Económicas} & \num{176855.8} & \num{160078.8} & \num{-16777.0} & \num{-24381.8} & \num{-13.2} \\
{07 Defensa Nacional} & \num{259433.8} & \num{158287.8} & \num{-101146.1} & \num{-112301.7} & \num{-41.5} \\
{09 Infraestructura, Comunicaciones y Transportes} & \num{85688.3} & \num{147511.5} & \num{61823.2} & \num{58138.6} & \num{65.1} \\
{18 Energía} & \num{167736.2} & \num{138307.4} & \num{-29428.9} & \num{-36641.5} & \num{-20.9} \\
{25 Previsiones y Aportaciones para los Sistemas de Educación Básica, Normal, Tecnológica y de Adultos} & \num{82518.5} & \num{81668.0} & \num{-850.5} & \num{-4398.8} & \num{-5.1} \\
{08 Agricultura y Desarrollo Rural} & \num{74109.6} & \num{74515.2} & \num{405.6} & \num{-2781.1} & \num{-3.6} \\
{03 Poder Judicial} & \num{78327.3} & \num{70983.6} & \num{-7343.7} & \num{-10711.7} & \num{-13.1} \\
{36 Seguridad y Protección Ciudadana} & \num{105838.8} & \num{70422.2} & \num{-35416.6} & \num{-39967.7} & \num{-36.2} \\
{12 Salud} & \num{96990.0} & \num{66693.2} & \num{-30296.8} & \num{-34467.4} & \num{-34.1} \\
{13 Marina} & \num{71888.2} & \num{65888.7} & \num{-5999.5} & \num{-9090.7} & \num{-12.1} \\
{34 Erogaciones para los Programas de Apoyo a Ahorradores y Deudores de la Banca} & \num{62489.4} & \num{52451.9} & \num{-10037.5} & \num{-12724.6} & \num{-19.5} \\
{30 Adeudos de Ejercicios Fiscales Anteriores} & \num{44050.6} & \num{45800.0} & \num{1749.4} & \num{-144.8} & \num{-0.3} \\
{16 Medio Ambiente y Recursos Naturales} & \num{70245.5} & \num{44370.5} & \num{-25875.0} & \num{-28895.5} & \num{-39.4} \\
{15 Desarrollo Agrario, Territorial y Urbano} & \num{12880.3} & \num{38048.0} & \num{25167.8} & \num{24613.9} & \num{183.2} \\
{38 Ciencia, Humanidades, Tecnología e Innovación} & \num{33170.7} & \num{33295.9} & \num{125.2} & \num{-1301.2} & \num{-3.8} \\
\bottomrule
\end{tabular}
\par\vspace{2pt}\footnotesize Fuente: elaboración propia del ITED con los analíticos del PEF aprobado 2024 y 2025. Las diferencias en millones de pesos se reportan dos veces, nominal y real; la real está en pesos de 2025 con el deflactor de 4.3 \% que declara el CGPE.
\end{table}


\marginal{la mitad de lo que se mueve no se mueve}
Los seis ramos que más caen en términos reales son Defensa Nacional (112,301.7
millones de pesos de 2025), Aportaciones Federales (46,865.5), Seguridad y
Protección Ciudadana (39,967.7), Energía (36,641.5), Salud (34,467.4) y Medio
Ambiente (28,895.5). Los que más suben son Deuda Pública (82,550.6),
Infraestructura, Comunicaciones y Transportes (58,138.6), Entidades no
Sectorizadas (29,725.9), Desarrollo Agrario (24,613.9), Participaciones (23,121.9)
y Bienestar (12,561.7).

\textbf{Dos de esos movimientos no son decisiones de gasto y hay que decirlo antes
de comentarlos.}

\subsection{Defensa no pierde 112 mil millones}

El Ramo 07 pasa de 259,433.8 a 158,287.8 millones de pesos. La caída entera vive
en dos unidades responsables que no son militares: \textbf{Tren Maya, S.A. de
C.V.} baja de 125,937.3 a 40,827.8, y el grupo aeroportuario y ferroviario de
15,172.8 a 2,275.5.

Y ese dinero no desapareció: se mudó. La función Transporte del Gobierno Federal
pasa de 212,619.2 a \textbf{214,635.9 millones de pesos}, prácticamente plana, con
una caída real de 3.2\%. Lo que cambió es quién la ejecuta: Defensa baja de
127,556.9 a 41,772.1, la Secretaría de Infraestructura sube de 80,811.1 a
147,181.7, y Marina de 4,251.2 a 25,682.1. Dentro de la Secretaría, la dirección
general de desarrollo ferroviario pasa de 9,120.8 a 92,028.2 millones, además de
cambiar de clave.

\textbf{La lectura correcta es que la obra ferroviaria vuelve a la administración
civil, no que el gasto militar se recorte.} La función Seguridad Nacional, que es
donde vive el gasto militar propiamente dicho, apenas se mueve: 103,530.1 a
101,948.0 millones en Defensa.

\subsection{Ciencia y tecnología: un ramo que cambia de nombre}

El Ramo 38 se llamaba ``Humanidades, Ciencias, Tecnologías e Innovación'' en 2024
y ``Ciencia, Humanidades, Tecnología e Innovación'' en 2025. Su monto pasa de
33,170.7 a 33,295.9 millones, es decir cae 4.1\% en términos reales.
\textbf{Cualquier comparación que agrupe por el nombre del ramo y no por su número
producirá un ramo que desaparece con 33 mil millones y otro que aparece con 33
mil, y una variación de 100\% que no existe.} Se anota porque es el error de
perímetro característico del género y porque la herramienta de esta casa lo
cometió antes de corregirse.

\section{Por finalidad y función}

\begin{table}[htbp]
\centering
\caption{Gasto del Gobierno Federal por finalidad y función (mdp)}
\label{tab:C4_4}
\footnotesize
\setlength{\tabcolsep}{4pt}
\begin{tabular}{L{3.4cm}SSSSS}
\toprule
{Concepto} & {2024} & {2025} & {Dif. nominal} & {Dif. real} & {Var. real \%} \\
\midrule
\num{2.6} & \num{1917554.8} & \num{2078757.6} & \num{161202.7} & \num{78747.9} & \num{3.9} \\
\num{4.2} & \num{1262789.1} & \num{1340210.9} & \num{77421.8} & \num{23121.9} & \num{1.8} \\
\num{4.1} & \num{1023011.5} & \num{1149551.6} & \num{126540.1} & \num{82550.6} & \num{7.7} \\
\num{2.5} & \num{1032621.4} & \num{1084590.9} & \num{51969.4} & \num{7566.7} & \num{0.7} \\
\num{2.2} & \num{333451.6} & \num{372916.8} & \num{39465.2} & \num{25126.8} & \num{7.2} \\
\num{2.3} & \num{439024.9} & \num{326133.5} & \num{-112891.4} & \num{-131769.5} & \num{-28.8} \\
\num{3.3} & \num{264603.7} & \num{242430.5} & \num{-22173.2} & \num{-33551.1} & \num{-12.2} \\
\num{3.5} & \num{212619.2} & \num{214635.9} & \num{2016.7} & \num{-7125.9} & \num{-3.2} \\
\num{1.6} & \num{150602.8} & \num{139011.9} & \num{-11590.9} & \num{-18066.8} & \num{-11.5} \\
\num{1.2} & \num{135515.6} & \num{128794.2} & \num{-6721.4} & \num{-12548.5} & \num{-8.9} \\
\num{3.2} & \num{90323.1} & \num{71802.6} & \num{-18520.6} & \num{-22404.5} & \num{-23.8} \\
\num{3.8} & \num{60385.7} & \num{60798.3} & \num{412.6} & \num{-2184.0} & \num{-3.5} \\
\num{1.5} & \num{81656.8} & \num{54494.1} & \num{-27162.7} & \num{-30674.0} & \num{-36.0} \\
\num{1.7} & \num{88200.2} & \num{53102.5} & \num{-35097.7} & \num{-38890.3} & \num{-42.3} \\
\num{4.3} & \num{62489.4} & \num{52451.9} & \num{-10037.5} & \num{-12724.6} & \num{-19.5} \\
\num{4.4} & \num{44050.6} & \num{45800.0} & \num{1749.4} & \num{-144.8} & \num{-0.3} \\
\num{1.3} & \num{45022.0} & \num{38738.9} & \num{-6283.1} & \num{-8219.1} & \num{-17.5} \\
\num{3.1} & \num{64048.4} & \num{32606.7} & \num{-31441.7} & \num{-34195.8} & \num{-51.2} \\
\bottomrule
\end{tabular}
\par\vspace{2pt}\footnotesize Fuente: elaboración propia del ITED con los analíticos del PEF aprobado 2024 y 2025. Las diferencias en millones de pesos se reportan dos veces, nominal y real; la real está en pesos de 2025 con el deflactor de 4.3 \% que declara el CGPE.
\end{table}


El movimiento funcional más grande del paquete es la caída de la función Salud,
que en el Gobierno Federal pasa de 439,024.9 a 326,133.5 millones de pesos.
\textbf{No se comenta aquí porque no significa lo que parece}; el capítulo 6 la
descompone y muestra que la mayor parte es un adeudo liquidado que no se repite y
un programa que cambió de ramo.

Los demás: el costo financiero sube 82,550.6 millones reales, la protección social
78,747.9, y bajan asuntos de orden público (38,890.3), asuntos económicos y
laborales (34,195.8), combustibles y energía (33,551.1) y asuntos financieros y
hacendarios (30,674.0).

\section{Lo que está comprometido}

El Anexo 3 del decreto publica dos cifras que este documento usa como columna
vertebral: las previsiones para gastos obligatorios, \textbf{6,046,446.8 millones
de pesos}, y las mismas con pensiones y jubilaciones, \textbf{7,684,111.9}.

De ahí sale un dato que conviene mirar despacio: \textbf{los gastos obligatorios
con pensiones equivalen al 82.6\% del gasto neto total} y al 145\% de la
recaudación tributaria. Y su diferencia, 1,637,665.1 millones, es exactamente el
agregado de pensiones y jubilaciones de la clasificación económica, lo que permite
adjudicar ese perímetro sin depender de los analíticos.

\marginal{una razón que se retira}
La razón que esta casa venía siguiendo para este capítulo, \emph{pensiones más
costo financiero sobre impuestos}, tenía un punto de 0.60 para 2023 que
\textbf{no se reproduce con ninguna de las dos definiciones de pensiones}: con el
perímetro amplio da 0.583 y con la clasificación económica 0.522. El punto se
retira y la serie arranca en 2024 con la definición declarada: \textbf{0.559 en
2024 y 0.571 en 2025}. Publicar una serie cuya base no cierra vale menos que
rehacerla.

\clearpage
